\chapter{Antecedentes}

\section{Definición de causalidad}
[Desarrollar las principales definiciones de causalidad.]

\section{La jerarquía causal de Pearl}
Pearl plantea como partida tres habilidades distintas de la cognición: ver, hacer e imaginar
o hipotetizar/anticipar(Rosen). \cite{pearl2018bookOfWhy}
\begin{itemize}
    \item \textbf{Ver - Imitar} Nos ayuda a percibir regularidades las observaciones. Estas
Rosen las pondría como la identificacion de perceptos. 
    \item \textbf{Hacer - Planear} Predecir los efectos provocados por cierta perturbacion en el ambiente
y tomar una decision para obtener el resultado deseado. Esto incluye la habilidad de
tomar herramientas intencionalmente para lograr el objetivo. 
    \item \textbf{Imaginar - Hipotetizar} Cuando tenemos una teoría de porque funciona una herramienta o cuando no, 
podemos hipotetizar y mejorar estas herramientas para futuros escenarios. 
\end{itemize}

Si una de las intenciones de Turing era intentar definir un agente coginitvo en cuanto a
las preguntas que puede responder. La intencion de Pearl aqui plantea, en cada uno de los 
niveles, distintas dificultades. Para Pearl un sistema que realmente presente cognicion 
deberia ser posible de reproducir estos tres tipos de inferencias causales. 

Los tres puntos anteriores representan entonces, la intuicion por detras de la 
formulacion matematica, de dichas relaciones. 

\begin{itemize}
    \item L1. \textbf{Asociativo}. ¿Que pasa si veo ? ¿Como cambiaria mi creencia en Y viendo X? En ejemplos 
    de la vida real sería que me dice un sintoma acerca de una enfermedad ? que me dice el resultado de una encuesta para predecir 
    el resultado final de las elecciones ? 
    \item L2. \textbf{Intervencion}. Que pasaria si hago \dots ? Como seria Y si hago X ? Como puedo hacer que pasae Y ?
    \item L3. \textbf{Contrafactual}. ¿Si tomo una aspirina, se me quitaria mi migraña ?¿Que pasa si prohibimos los cigarros?
\end{itemize}

¿Que es lo que puede hacer un razonador causal?¿Que puede computar un organismo con un modelo causal, 
a falta de el, que le permite predecir? 

Aqui Rosen diferia, en cuanto a que el proceso de inferencia de un organismo, 
no de depende solo de su estructura causal sino tambien viene desde una idea que 
engloba el completo funcionamiento del organismo, definiendo la coginicion y agencia del 
organismo no solo en la estructura causal del fenomeno que se esta tratando de predecir. 

Uno de los objetivos principales de la version fuerte de IA, es reproducir la inteligencia humana. En lugar
de tener una inteligencia genuina, nos encontramos con habilidades impresionantes. La diferencia recae en la falta 
de un modelo de la realidad. 

Muchos de los sistamas de IA se basan especialmente en la asociacion. Toman una colección gigante de observaciones para 
encajar una funcion de la misma manera que un estadista podria querer encajar una linea para pasar por una 
coleccion de puntos. Esto pone una barrera, para que los sistemas puedan ser capaces de reaccionar a circunstancias fuera del 
dataset. Como en los coches de manejo automatico, ¿que pasa si nos encontramos por accidente una vaca en la carretera, un bache
o un borracho ?¿Como podemos hacer un deduccion generalizada? Esta falta de flexibilidad y adaptacion es ineblitable 
en la primerca jerarquia causal. 


\section{Modelo causal estructural}
[Describir los modelos causales estructurales y los grafos causales.]

\section{Relación de modelado de Rosen}
[Presentar la relación de modelado y su vínculo con los sistemas anticipatorios.]

\section{Integración teórica}
[Proponer el criterio para distinguir inferencia causal de simulación predictiva.]
